\begin{spacing}{1}
\noindent Os Regimes Próprios de Previdência Social (RPPS) são os sistemas responsáveis por prover benefícios pós-emprego aos servidores públicos efetivos. Nesse contexto, a Portaria MTP nº 1.467/2022 estabelece parâmetros e diretrizes para estas organizações. Em seu Art. 36, a referida Portaria estabelece como critério mínimo para a seleção de tábuas de mortalidade a exigência de que a tábua adotada apresente, na idade média da massa segurada, expectativa de vida igual ou superior à da tábua de referência (IBGE anual extrapolada por sexo). Este trabalho analisa se o critério de seleção de tábuas de mortalidade estabelecido pelo Art. 36 da referida Portaria é suficiente para garantir que a provisão matemática calculada com a tábua adotada seja igual ou superior àquela obtida com a tábua de referência. Para isso, foram realizadas simulações com quatro perfis etários distintos — jovem, equilibrado, idoso e bimodal —, cada um composto por 100.000 segurados de cada sexo, considerando um cenário simplificado com único benefício de aposentadoria por idade, salário inicial uniforme e ausência de decrementos por invalidez ou desligamento. As provisões foram calculadas pelos métodos INE, CUP e PNI, aplicando-se o critério de seleção de tábuas conforme o fluxo normativo do Art. 36. Os resultados evidenciaram inconsistências em todos os perfis etários, com maior incidência em massas jovens, especialmente para o sexo masculino. A investigação revelou que a assimetria entre os níveis de sobrevivência nas faixas contributivas (18–64 anos) e nas faixas em gozo de benefício (65 anos em diante) é o mecanismo causador. Quando a tábua comparativa projeta maior sobrevivência nas idades contributivas do que nas aposentadoras, a variação no valor das contribuições supera a variação no valor dos benefícios, resultando em provisão matemática inferior. Para auxiliar o diagnóstico prévio, propõe-se o índice $I_F$, que quantifica o diferencial de sobrevivência ponderado pela estrutura etária da massa. Conclui-se que o critério do Art. 36, ao utilizar exclusivamente a expectativa de vida na idade média como parâmetro comparativo, pode ser insuficiente para garantir que a tábua adotada produza uma provisão matemática igual ou superior à da tábua de referência em todas as configurações demográficas possíveis. 
\end{spacing}
\vspace{1.0cm}

\noindent \textbf{Palavras-chave}: Regimes Próprios de Previdência Social. Tábuas de mortalidade. Provisões matemáticas. Portaria MTP nº 1.467/2022. Avaliação atuarial.